% Algebra 2 Lesson 5-5 fast-track packet: Function Operations (enVision A2 2024 TE, Example 5 and Example 6 only)
% Kept: Example 5 (quiz items 1, 4, 5) and Example 6 (the order of a composition; it carries the lesson's only DOK 3 items, 16 and 32).
% Cut: Examples 1-4 (quiz items 2 and 3 are left off the short form). DOK 3 = Practice 16, the Exit question; Practice 32 is the extension.
% Shared preamble for the "fast" (hacked-down) Algebra 2 lesson packets, Oct 2026.
% Each packet: % Shared preamble for the "fast" (hacked-down) Algebra 2 lesson packets, Oct 2026.
% Each packet: % Shared preamble for the "fast" (hacked-down) Algebra 2 lesson packets, Oct 2026.
% Each packet: \input{fast-packet-template} then \begin{document} ... \end{document}.
\documentclass[12pt]{article}
\usepackage[margin=0.7in]{geometry}
\usepackage{amsmath,amssymb}
\usepackage{booktabs,tabularx,array}
\usepackage{enumitem}
\usepackage[T1]{fontenc}
\usepackage[utf8]{inputenc}
\usepackage{lmodern}
\usepackage{tikz}
\usepackage{pgfplots}
\pgfplotsset{compat=1.18}
\usepackage{xcolor}
\setlist{nosep,leftmargin=1.4em}
\setlength{\parindent}{0pt}
\setlength{\parskip}{6pt}
\renewcommand{\arraystretch}{1.35}
\newcolumntype{Y}{>{\raggedright\arraybackslash}X}
\pagestyle{empty}

% blanks and workspace
\newcommand{\blank}[1][2.6cm]{\rule[-2pt]{#1}{0.5pt}}
\newcommand{\lines}[1][1]{\par\vspace{2pt}\foreach \i in {1,...,#1}{\rule{0pt}{22pt}\rule{\linewidth}{0.4pt}\par}}
\newcommand{\wspart}[2]{\vspace{4pt}{\large\bfseries #1: #2}\par\vspace{-6pt}\rule{\textwidth}{0.4pt}\par}

% header: \packethead{Lesson 1-3}{Piecewise-Defined Functions}
\newcommand{\packethead}[2]{%
{\LARGE\bfseries #2}\hfill Name \blank[5cm]\\[2pt]
Algebra 2 \quad #1 (fast track) \hfill Date \blank[2.5cm] \quad Period \blank[1.2cm]\par\vspace{2pt}}

% boxed goals: \goals{learning goal}{language goal}
\newcommand{\goals}[2]{%
\fbox{\parbox{\dimexpr\textwidth-2\fboxsep-2\fboxrule}{%
\textbf{Learning goal:} #1\\[2pt]
\textbf{Language goal:} #2}}\par}

% quiz target box: \quiztarget{text}
\newcommand{\quiztarget}[1]{%
\vspace{4pt}\colorbox{gray!12}{\parbox{\dimexpr\textwidth-2\fboxsep}{\textbf{On the quiz:} #1}}\par}

% IXL table row helper: code & skill & which quiz question it feeds
\newcommand{\ixlrow}[3]{\texttt{#1} & #2 & #3 \\}

% one grid style for every graph; override xmin/xmax/... per plot
\pgfplotsset{sheetgrid/.style={
  axis lines=middle, grid=both, grid style={gray!45},
  xmin=-10.5, xmax=10.5, ymin=-10.5, ymax=10.5,
  xtick={-10,-8,...,10}, ytick={-10,-8,...,10},
  x=0.28cm, y=0.28cm,
  ticklabel style={font=\scriptsize, fill=white, inner sep=1pt},
  xlabel={$x$}, ylabel={$y$},
  axis line style={thick}, clip=true}}
% unit grid for small windows
\pgfplotsset{unitgrid/.style={sheetgrid,
  xmin=-5.5, xmax=5.5, ymin=-5.5, ymax=5.5,
  xtick={-5,...,5}, ytick={-5,...,5}, x=0.5cm, y=0.5cm}}
% open and closed endpoints
\tikzset{closed/.style={circle, fill, inner sep=1.6pt},
         open/.style={circle, draw, fill=white, inner sep=1.4pt}}
 then \begin{document} ... \end{document}.
\documentclass[12pt]{article}
\usepackage[margin=0.7in]{geometry}
\usepackage{amsmath,amssymb}
\usepackage{booktabs,tabularx,array}
\usepackage{enumitem}
\usepackage[T1]{fontenc}
\usepackage[utf8]{inputenc}
\usepackage{lmodern}
\usepackage{tikz}
\usepackage{pgfplots}
\pgfplotsset{compat=1.18}
\usepackage{xcolor}
\setlist{nosep,leftmargin=1.4em}
\setlength{\parindent}{0pt}
\setlength{\parskip}{6pt}
\renewcommand{\arraystretch}{1.35}
\newcolumntype{Y}{>{\raggedright\arraybackslash}X}
\pagestyle{empty}

% blanks and workspace
\newcommand{\blank}[1][2.6cm]{\rule[-2pt]{#1}{0.5pt}}
\newcommand{\lines}[1][1]{\par\vspace{2pt}\foreach \i in {1,...,#1}{\rule{0pt}{22pt}\rule{\linewidth}{0.4pt}\par}}
\newcommand{\wspart}[2]{\vspace{4pt}{\large\bfseries #1: #2}\par\vspace{-6pt}\rule{\textwidth}{0.4pt}\par}

% header: \packethead{Lesson 1-3}{Piecewise-Defined Functions}
\newcommand{\packethead}[2]{%
{\LARGE\bfseries #2}\hfill Name \blank[5cm]\\[2pt]
Algebra 2 \quad #1 (fast track) \hfill Date \blank[2.5cm] \quad Period \blank[1.2cm]\par\vspace{2pt}}

% boxed goals: \goals{learning goal}{language goal}
\newcommand{\goals}[2]{%
\fbox{\parbox{\dimexpr\textwidth-2\fboxsep-2\fboxrule}{%
\textbf{Learning goal:} #1\\[2pt]
\textbf{Language goal:} #2}}\par}

% quiz target box: \quiztarget{text}
\newcommand{\quiztarget}[1]{%
\vspace{4pt}\colorbox{gray!12}{\parbox{\dimexpr\textwidth-2\fboxsep}{\textbf{On the quiz:} #1}}\par}

% IXL table row helper: code & skill & which quiz question it feeds
\newcommand{\ixlrow}[3]{\texttt{#1} & #2 & #3 \\}

% one grid style for every graph; override xmin/xmax/... per plot
\pgfplotsset{sheetgrid/.style={
  axis lines=middle, grid=both, grid style={gray!45},
  xmin=-10.5, xmax=10.5, ymin=-10.5, ymax=10.5,
  xtick={-10,-8,...,10}, ytick={-10,-8,...,10},
  x=0.28cm, y=0.28cm,
  ticklabel style={font=\scriptsize, fill=white, inner sep=1pt},
  xlabel={$x$}, ylabel={$y$},
  axis line style={thick}, clip=true}}
% unit grid for small windows
\pgfplotsset{unitgrid/.style={sheetgrid,
  xmin=-5.5, xmax=5.5, ymin=-5.5, ymax=5.5,
  xtick={-5,...,5}, ytick={-5,...,5}, x=0.5cm, y=0.5cm}}
% open and closed endpoints
\tikzset{closed/.style={circle, fill, inner sep=1.6pt},
         open/.style={circle, draw, fill=white, inner sep=1.4pt}}
 then \begin{document} ... \end{document}.
\documentclass[12pt]{article}
\usepackage[margin=0.7in]{geometry}
\usepackage{amsmath,amssymb}
\usepackage{booktabs,tabularx,array}
\usepackage{enumitem}
\usepackage[T1]{fontenc}
\usepackage[utf8]{inputenc}
\usepackage{lmodern}
\usepackage{tikz}
\usepackage{pgfplots}
\pgfplotsset{compat=1.18}
\usepackage{xcolor}
\setlist{nosep,leftmargin=1.4em}
\setlength{\parindent}{0pt}
\setlength{\parskip}{6pt}
\renewcommand{\arraystretch}{1.35}
\newcolumntype{Y}{>{\raggedright\arraybackslash}X}
\pagestyle{empty}

% blanks and workspace
\newcommand{\blank}[1][2.6cm]{\rule[-2pt]{#1}{0.5pt}}
\newcommand{\lines}[1][1]{\par\vspace{2pt}\foreach \i in {1,...,#1}{\rule{0pt}{22pt}\rule{\linewidth}{0.4pt}\par}}
\newcommand{\wspart}[2]{\vspace{4pt}{\large\bfseries #1: #2}\par\vspace{-6pt}\rule{\textwidth}{0.4pt}\par}

% header: \packethead{Lesson 1-3}{Piecewise-Defined Functions}
\newcommand{\packethead}[2]{%
{\LARGE\bfseries #2}\hfill Name \blank[5cm]\\[2pt]
Algebra 2 \quad #1 (fast track) \hfill Date \blank[2.5cm] \quad Period \blank[1.2cm]\par\vspace{2pt}}

% boxed goals: \goals{learning goal}{language goal}
\newcommand{\goals}[2]{%
\fbox{\parbox{\dimexpr\textwidth-2\fboxsep-2\fboxrule}{%
\textbf{Learning goal:} #1\\[2pt]
\textbf{Language goal:} #2}}\par}

% quiz target box: \quiztarget{text}
\newcommand{\quiztarget}[1]{%
\vspace{4pt}\colorbox{gray!12}{\parbox{\dimexpr\textwidth-2\fboxsep}{\textbf{On the quiz:} #1}}\par}

% IXL table row helper: code & skill & which quiz question it feeds
\newcommand{\ixlrow}[3]{\texttt{#1} & #2 & #3 \\}

% one grid style for every graph; override xmin/xmax/... per plot
\pgfplotsset{sheetgrid/.style={
  axis lines=middle, grid=both, grid style={gray!45},
  xmin=-10.5, xmax=10.5, ymin=-10.5, ymax=10.5,
  xtick={-10,-8,...,10}, ytick={-10,-8,...,10},
  x=0.28cm, y=0.28cm,
  ticklabel style={font=\scriptsize, fill=white, inner sep=1pt},
  xlabel={$x$}, ylabel={$y$},
  axis line style={thick}, clip=true}}
% unit grid for small windows
\pgfplotsset{unitgrid/.style={sheetgrid,
  xmin=-5.5, xmax=5.5, ymin=-5.5, ymax=5.5,
  xtick={-5,...,5}, ytick={-5,...,5}, x=0.5cm, y=0.5cm}}
% open and closed endpoints
\tikzset{closed/.style={circle, fill, inner sep=1.6pt},
         open/.style={circle, draw, fill=white, inner sep=1.4pt}}

\begin{document}

\packethead{Lesson 5-5}{Function Operations}
\goals{I can write the rule and state the domain of a composite function, and use the order of a composition to answer a real-world question.}{I can explain how I composed two functions using the words composition, inside function, substitute, order and domain.}

\quiztarget{(1) $g(f(x))+g(x)$: \emph{substitute} the whole rule for $f$ into $g$ (Example 5B, Try It 5), then add $g(x)$; the wrong \emph{order}, $f(g(x))$, is a choice. (2) Revenue and (3) a quotient $\frac fg$ and its \emph{domain} (full Savvas quiz only). (4, bonus) Work inside out with square roots: $k(3)$ first, then $h$ of that number (Example 5A, the Exit question). (5) The rule and \emph{domain} of $f\circ g$ when the inside function $g$ has the square root (Example 5A, checks 1 and 2).}

\vspace{2pt}
\fbox{\parbox{\dimexpr\textwidth-2\fboxsep-2\fboxrule}{%
\textbf{Composite function.} Apply the rule for $g$, then apply the rule for $f$ to the result:
\[(f\circ g)(x)=f(g(x))\qquad\text{$g$ is the \textbf{inside function} (first); $f$ is the outside function (second).}\]
The operation $\circ$ is called \textbf{composition of functions}. The \textbf{domain} of $f\circ g$ is every $x$ in the domain of $g$ (check 1) for which $g(x)$ is in the domain of $f$ (check 2).

\textbf{Order matters:} for $f(x)=3x+5$, $g(x)=x-3$: \ \mbox{$f\circ g=3(x-3)+5=3x-4$},\newline but \mbox{$g\circ f=(3x+5)-3=3x+2$}.}}

\wspart{Example 5}{Write a Rule for a Composite Function}
What is the rule for the composition $(f\circ g)(x)$? \emph{Substitute} the rule for $g$ wherever $x$ appears in $f$.

\begin{minipage}[t]{0.47\textwidth}
\textbf{A.} $f(x)=\sqrt{x+7}$ \ and \ $g(x)=2x-5$
\begin{align*}
(f\circ g)(x)&=f(g(x))\\
&=f(2x-5)\\
&=\sqrt{(2x-5)+7}\\
&=\sqrt{2x+2}
\end{align*}
\textbf{Domain.} $g$ takes all real numbers; under the root, $2x+2\geq0$, so the domain of $f\circ g$ is $x\geq-1$.

\textbf{Check} $x=1$: $g(1)=-3$, $f(-3)=\sqrt{4}=2$; and $\sqrt{2(1)+2}=2$. \checkmark
\end{minipage}\hfill
\begin{minipage}[t]{0.49\textwidth}
\textbf{B.} $f(x)=x^2+x+2$ \ and \ $g(x)=4-x$
\begin{align*}
(f\circ g)(x)&=f(g(x))\\
&=f(4-x)\\
&=(4-x)^2+(4-x)+2\\
&=16-8x+x^2+4-x+2\\
&=x^2-9x+22
\end{align*}
\textbf{Domain.} All real numbers.

\textbf{Watch:} \emph{every} $x$ in $f$ is replaced, including the lone $x$: it becomes $(4-x)$.

\textbf{Domain checks:} in A, check 2 sets the limit; when $g$ has the square root, check 1 does.
\end{minipage}

\vspace{4pt}
\textbf{Say it:} To find $(f\circ g)(x)$ in Part A, I \blank[2.4cm] the inside function \blank[2cm] for $x$ in $f$. The domain is $x\geq-1$ because \dots\lines[2]

\newpage

\wspart{Try It 5}{Write the rules for two compositions}
Identify the rules for $f\circ g$ and $g\circ f$.

Hint: Write $f(g(x))$ and $g(f(x))$ first, then substitute and simplify.

\begin{minipage}[t]{0.48\textwidth}
\textbf{a.} $f(x)=x^3$, \ $g(x)=x+1$

\vspace{4pt}
$(f\circ g)(x)=f(g(x))=$\lines[6]

\vspace{6pt}
$(g\circ f)(x)=g(f(x))=$\lines[5]
\end{minipage}\hfill
\begin{minipage}[t]{0.48\textwidth}
\textbf{b.} $f(x)=3\sqrt{2x}$, \ $g(x)=4x+11$

\vspace{4pt}
$(f\circ g)(x)=f(g(x))=$\lines[6]

\vspace{6pt}
$(g\circ f)(x)=g(f(x))=$\lines[5]
\end{minipage}

\vspace{8pt}
\textbf{Say it:} In part a, $f\circ g$ and $g\circ f$ are \blank[2.2cm] (the same / different), so the \blank[2cm] of a composition \blank[2.2cm] (matters / does not matter), because \dots\lines[3]

\newpage

\wspart{Example 6}{Use a Composite Function Model}
A clothing store posts two discounts on social media. Followers may use both; they tell the cashier in which order to apply them. In which order should they ask for the discounts to save the most money?

\noindent\fbox{\parbox{7.6cm}{\centering\scriptsize EXCLUSIVELY FOR OUR SOCIAL MEDIA FOLLOWERS!\\ \normalsize\textbf{\$5 OFF} NEXT PURCHASE}}\hfill
\fbox{\parbox{7.6cm}{\centering\scriptsize EXCLUSIVELY FOR OUR SOCIAL MEDIA FOLLOWERS!\\ \normalsize\textbf{10\% OFF} NEXT PURCHASE}}\par

Let $x$ be the price of the purchase. Model each discount with a function:\\
\hspace*{1.5em}$f(x)=x-5$ \ (\$5 off) \qquad $g(x)=x-0.1x=0.9x$ \ (price minus 10\% of the price)

\begin{minipage}[t]{0.48\textwidth}
\textbf{10\% off, then \$5 off:} the rule for $f\circ g$ ($g$ first).
\begin{align*}
(f\circ g)(x)&=f(g(x))\\
&=f(0.9x)\\
&=0.9x-5
\end{align*}
\end{minipage}\hfill
\begin{minipage}[t]{0.48\textwidth}
\textbf{\$5 off, then 10\% off:} the rule for $g\circ f$ ($f$ first).
\begin{align*}
(g\circ f)(x)&=g(f(x))\\
&=g(x-5)\\
&=0.9(x-5)=0.9x-4.5
\end{align*}
\end{minipage}

\textbf{Test a price.} A shirt costs \$30.00 before the discounts.\\
$(f\circ g)(30)=0.9(30)-5=22$, or \$22.00. \qquad $(g\circ f)(30)=0.9(30)-4.5=22.5$, or \$22.50.

\textbf{Answer:} ask for \textbf{10\% off first, then \$5 off}. The 10\% is taken off the larger price, so it saves more.

\textbf{Say it:} The order ``10\% first, then \$5'' is the composition \blank[1.6cm], because $g$ is the \blank[2.4cm]. It is the better deal because \dots\lines[1]

\wspart{Try It 6}{Model a discount and a sales tax}
As a member of the Games Shop rewards program, you get a 20\% discount on purchases. All sales are subject to a 6\% sales tax. Write functions to model the discount and the sales tax, then identify the rule for the composition function that calculates the final price you would pay at Games Shop.

\begin{minipage}[t]{0.48\textwidth}
Let $x=$ \blank[4.2cm]

\vspace{6pt}
Discount: \ $d(x)=\blank[3.2cm]$

\vspace{6pt}
Sales tax: \ $t(x)=\blank[3.2cm]$
\end{minipage}\hfill
\begin{minipage}[t]{0.48\textwidth}
\raggedright Final price (which function is the inside function?):

\vspace{6pt}
$(\blank[0.8cm]\circ\blank[0.8cm])(x)=\blank[4cm]$
\lines[2]
\end{minipage}

\textbf{Say it:} In my composition the inside function is \blank[1.2cm], so the \blank[2.2cm] is applied first. Switching the order here \blank[3cm] (changes / does not change) the final price, because \dots\lines[1]

\newpage

\wspart{IXL}{Practice for the quiz}
Do the starred skills first. \textbf{SmartScore 70 or above on each starred ($\bigstar$) skill replaces your quiz score.}

\vspace{2pt}
\begin{tabularx}{\textwidth}{|l|Y|Y|}
\hline
\textbf{Code} & \textbf{Skill} & \textbf{Feeds} \\ \hline
\ixlrow{EKJ $\bigstar$}{Composition of linear and quadratic functions: find an equation}{Quiz 1: $g(f(x))$ with a quadratic inside; Example 5B, Try It 5a}
\ixlrow{RSP $\bigstar$}{Composition of linear functions: find an equation}{Quiz 1 and 5: the rule for $f\circ g$ in the right order; Example 6, Try It 6}
\ixlrow{P9T $\bigstar$}{Composition of linear and quadratic functions: find a value}{Quiz 4 (bonus): inside out, one number at a time}
\ixlrow{QQD $\bigstar$}{Add and subtract functions}{Quiz 1: the $+\,g(x)$, by combining like terms}
\ixlrow{PS2}{Evaluate functions}{Warm-up, optional: substitute a number or an expression}
\hline
\end{tabularx}

\wspart{Exit question}{Compose two functions given as ordered pairs}
\fbox{\parbox{\dimexpr\textwidth-2\fboxsep-2\fboxrule}{%
Suppose two functions, $f(x)$ and $g(x)$, are only defined by the ordered pairs listed below.
\[f=(6,7),\ (5,2),\ (4,1),\ (10,8)\qquad\qquad g=(5,4),\ (3,6),\ (1,5),\ (2,10)\]
Find the ordered pairs that comprise $(f\circ g)(x)$.

\vspace{4pt}
\begin{minipage}[t]{0.44\textwidth}
\vspace{0pt}
\renewcommand{\arraystretch}{1.6}
\begin{tabular}{|c|c|c|}
\hline
$x$ & $g(x)$ (inside) & $f(g(x))$ (outside) \\ \hline
\phantom{00} & & \\ \hline
& & \\ \hline
& & \\ \hline
& & \\ \hline
\end{tabular}
\end{minipage}\hfill
\begin{minipage}[t]{0.52\textwidth}
\vspace{0pt}
$(f\circ g)(x)=\{\,\blank[5.2cm]\,\}$

\vspace{6pt}
\textbf{Say it:} The domain of $f\circ g$ is $\{\blank[2.6cm]\}$ because each output of the inside function $g$ is \dots\lines[3]
\end{minipage}
}}

\vspace{6pt}
\textbf{Before the quiz, finish these sentences.}

To find $(f\circ g)(x)$, I substitute the inside function \blank[2.2cm] for $x$ in \blank[2.2cm].

The domain of $f\circ g$ can be smaller than the domain of $g$ because \blank[6cm].

In a composition the order matters: $f\circ g$ applies \blank[1.2cm] first, and $g\circ f$ applies \blank[1.2cm] first.

\vspace{4pt}
{\small\textbf{Extension (self-paced, DOK 3, Practice 32):} A music store takes \$5 off a purchase of \$20 or more, and you save an additional 15\% off your total purchase when you open a charge account. Use composition of functions to find the sale price of a \$90 purchase when the \$5 off is applied first, then when the 15\% off is applied first. In which order is the deal better for the consumer? Explain.}

\end{document}
